% =============================================================================
\section{Binding Modules}
\label{sec:modules}
% =============================================================================
We describe the bindings generated by our system for various modules
and examplify the use of the generated bindings.
%
When developing code in Python that uses the bindings, the statement
that imports the binding library, e.g., \pyLstinline{import CGALPY}
must precede any statement that use the binding. This statement is
omitted in all examples hereafter.

% -----------------------------------------------------------------------------
\subsection{Kernel Bindings}
\label{ssec:modules:kernel}
% -----------------------------------------------------------------------------
The \iiDandiiiDGeometryKernelPackage{}
package~\cite{cgal:bfghhkps-lgk23-21a} consists of constant-size
non-modifiable geometric primitive objects and operations on these
objects. The objects are sets of points in $d$-dimensional affine
Euclidean space, where $d = 2,3$.  Each point is uniquely represented
either by Cartesian coordinates or by homogeneous coordinates. An
object type can be defined either precatively as a member of a kernel
type or imperatively as a global class-template parameterized by a
kernel type, defined in the C++ \cppLstinline{CGAL} namespace. For
example, assume that \cppLstinline{Kernel} is a kernel type; the type
that represents a two-dimensional point of this kernel is either
\cppLstinline{Kernel::Point_2} or
\cppLstinline{CGAL::Point_2<Kernel>}.  We generate bindings only for
the latter types; that is, when defining the two-dimensional point in
Python code amounts to the code below.\footnote{When writing code in
  C++ the precative style is advantageous, as it enables the extension of
  kernel object types, which is irrelevant when generating bindings.}
\begin{lstlisting}[style=pythonStyle]
>>> Ker = CGALPY. Ker     # define the module namespace
>>> Point_2 = Ker.Point_2 # define the point type in the module namespace
>>> p = Point_2(0,0);
\end{lstlisting}
The kernel type determines the underlying number types used for
calculations. Table~\ref{tab:ker} lists the \cmake{} flags associated
with the \iiDandiiiDGeometryKernelPackage{} module.  The kernel type
is an instance of a chain of C++ class templates. For convenience,
\cgal{} provides the following predefned types of generally useful
kernels:
\begin{compactenum}
\item \cppLstinline{Exact_predicates_inexact_constructions_kernel}---provides
  exact geometric predicates, but geometric constructions may be inexact
  due to round-off errors.
  %
\item \cppLstinline{Exact_predicates_exact_constructions_kernel}---provides
  exact geometric constructions, in addition to exact geometric predicates.
  %
\item \cppLstinline{Exact_predicates_exact_constructions_kernel_with_sqrt}---same
  as\\
  \cppLstinline{Exact_predicates_exact_constructions_kernel}, but the
  number type is a model of the concept that requires operations that perform
  square roots, namely \cppLstinline{FieldWithSqrt}.\footnote{See \url{https://doc.cgal.org/latest/Algebraic_foundations/classFieldWithSqrt.html}.}
  %
\item \cppLstinline{Exact_predicates_exact_constructions_kernel_with_kth_root}---same
  as \\
  \cppLstinline{Exact_predicates_exact_constructions_kernel}, but the number
  type is a model of the concept that requires operations that perform $k$-th
  roots, namely
  \cppLstinline{FieldWithKthRoot}.\footnote{See \url{https://doc.cgal.org/latest/Algebraic_foundations/classFieldWithKthRoot.html}}
  %
\item \cppLstinline{Exact_predicates_exact_constructions_kernel_with_root_of}---same\\
  as \cppLstinline{Exact_predicates_exact_constructions_kernel}, but the number
  type is a model of the concept that requires operations that computes the root of univariate polynomial, namely \cppLstinline{FieldWithRootOf}.\footnote{See \url{https://doc.cgal.org/latest/Algebraic_foundations/classFieldWithRootOf.html}.}
\end{compactenum}

%%%%%%%% Kernel Module Flags
\begin{table}[!htp]
  \caption{\iiDandiiiDGeometryKernelPackage{} module flags}
  \centering\begin{tikzpicture}
  \matrix (first) [tableModuleFlags,
    row 3/.style={nodes={minimum height=2.8em}}] {
Name & Type & Default & Description\\
\myLstinline{KERNEL\_NAME} & String & \myLstinline{epic} & The kernel type used\\
\myLstinline{KERNEL\_INTERSECTION\_BINDINGS} & Boolean & \myLstinline{true} & Determines whether to generate bindings for intersections\\
  };
 \end{tikzpicture}
 \label{tab:ker}
\end{table}

All the predefined types are of Cartesian
kernels. The \cmake{} flag \myLstinline{KERNEL\_NAME} specifies which
kernel type should be used for the generated bindings.  Currently,
only a limitted predefined types are supported; see Table~\ref{tab:ker:name}.

%%%%%%%% Kernel Name Options
\begin{table}[!htp]
 \caption{Kernel name options}
 \centering\begin{tikzpicture}
   \matrix (first) [tableName,
     column 1/.style={nodes={text width=8em}},
     column 2/.style={nodes={text width=41em}}
   ] {
   \myLstinline{KERNEL\_NAME} & Predefined Type\\
   \myLstinline{epic} &
     \cppLstinline{Exact_predicates_inexact_constructions_kernel}\\
   \myLstinline{epec} &
     \cppLstinline{Exact_predicates_exact_constructions_kernel}\\
   \myLstinline{epecws} &
     \cppLstinline{Exact_predicates_exact_constructions_kernel_with_sqrt}\\
  };
 \end{tikzpicture}
 \label{tab:ker:name}
\end{table}

The code excerpt below, in the \cmake{} language, sets the \cmake{}
flags for our first example. When \myLstinline{cmake} is applied with
these settings followed \myLstinline{cmake} a library called
\myLstinline{CGALPY}, which consists of the necessary bindings, is
generated.
\framedInputCmake[\normalsize]{../../cmake/tests/epec_fixed_release.cmake}

The following example computes the intersection point of two segments.
\framedInputPython[\normalsize][15]{../../src/python_scripts/cgalpy_examples/kernel_intersection.py}

% -----------------------------------------------------------------------------
\subsection{2D Arrangement Bindings}
\label{ssec:modules:aos2}
% -----------------------------------------------------------------------------
Given a surface $S$ in \Rthree{} and a set $\calC$ of curves embedded
in this surface, the curves subdivide $S$ into cells of dimension~$0$
(vertices),~$1$ (edges), and~$2$ (faces). This subdivision is the
arrangement $\calA(\calC)$ induced by $\calC$ on $S$. Arrangements
embedded in curved surfaces in \Rthree{} are generalizations of
arrangements embedded in the plane. The \iiDArrangementsPackage{}
package~\cite{cgal:wfzh-a2-21a} can be used to construct, maintain,
alter, and display 2D arrangements embedded in ruled curved
surfaces. It also supports queries on such arrangements, such as point
location, and other operations. One of the main components of
the \iiDArrangementsPackage{} package is the
\cppLstinline{Arrangement_2<Traits,Dcel>} class template. An instance of this
template is used to represent an arrangement embedded in the plane.
Table~\ref{tab:aos} lists the \cmake{} flags associated with
the \iiDArrangementsPackage{} module.  Currently, only instances of
the \cppLstinline{Arrangement_2<Traits,Dcel>} class template are
supported. A description of the two template parameters of this class
template follows.

%%%%%%%% Arrangement Module Flags
\begin{table}[!htp]
 \caption{\iiDArrangementsPackage{} module flags}
  \centering\begin{tikzpicture}
  \matrix (first) [tableModuleFlags,
    row 4/.style={nodes={minimum height=4.0em}}
  ] {
    Name & Type & Default & Description\\
    \myLstinline{AOS2\_GEOMETRY\_TRAITS\_NAME}    & String  &
      \myLstinline{segment}     & The geometry traits of a 2D arrangement\\
    \myLstinline{AOS2\_DCEL\_NAME}                & String  &
      \myLstinline{allExtended} & The DCEL of a 2D arrangement\\
    \myLstinline{AOS2\_POINT\_LOCATION\_BINDINGS} & Boolean &
      \myLstinline{true}        & Determines whether to generate bindings for point location and vertical ray shooting queries\\
  };
 \end{tikzpicture}
 \label{tab:aos}
\end{table}

\begin{itemize}
  \item The Traits template-parameter determines the family of planar
  curves that induce the arrangement. The parameter should be
  substituted with a model of the basic arrangement traits concept,
  and optionally additional geometry traits concepts. A model of the
  basic traits concept defines the types of x-monotone curves and
  two-dimensional points and supports basic geometric predicates on
  them. The list of supported traits class templates follows. For each
  class template we describe the family of curves it handles.
  %
  \begin{compactenum}
  \item \cppLstinline{Arr_non_caching_segment_basic_traits_2<Kernel>}---handles
    segments.
  \item \cppLstinline{Arr_segment_traits_2<Kernel>}---handles segments,
    where each segment is represented a by its supporting line in
    addition to its two endpoints.
  \item \cppLstinline{Arr_linear_traits_2<Kernel>}---handles linear curves,
    i.e., segments, rays, and lines.
  \item \cppLstinline{Arr_polyline_traits_2<SubcurveTraits_2>}
  \item \cppLstinline{Arr_polycurve_traits_2<SubcurveTraits_2>}
  \item \cppLstinline{Arr_circle_segment_traits_2<Kernel>}
  \item \cppLstinline{Arr_conic_traits_2<RatKernel, AlgKernel, NtTraits>}
  \item \cppLstinline{Arr_rational_function_traits_2<AlgebraicKernel_d_1>}
  \item \cppLstinline{Arr_Bezier_curve_traits_2<RatKernel, AlgKernel, NtTraits>}
  \item \cppLstinline{Arr_algebraic_segment_traits_2<Coefficient>}
  \end{compactenum}

  The \cmake{} flag \myLstinline{AOS2\_GEOMETRY\_TRAITS\_NAME}
  specifies which geometry traits should be used for the generated
  bindings; see Table~\ref{tab:aos:gt}.
  %
\item The Dcel template-parameter should be substituted with a class that
  models the \cppLstinline{ArrangementDcel} concept, which is used to
  represent the topological layout of the
  arrangement.\footnote{See \url{https://doc.cgal.org/latest/Arrangement_on_surface_2/classArrangementDcel.html}.}
  This parameter is substituted
  with \cppLstinline{Arr_default_dcel<Traits>} by default. In many
  applications it is necessary to extend the DCEL features. This is
  done by substituting the Dcel parameter with a different type, which can
  be an instance of one of the following class templates:
  \begin{compactenum}
  \item \cppLstinline{Arr_extended_vertex<VertexBase, VertexData>}---extends
    the vertex topological-features of the Dcel with a record of type
    \cppLstinline{VertexData}.
    %
  \item \cppLstinline{Arr_extended_halfedge<HalfedgeBase, HalfedgeData>}---extends
    the halfedge topological-fea-tures of the Dcel with a record of type
    \cppLstinline{HalfedgeData}.
    %
  \item \cppLstinline{Arr_extended_face<FaceBase, FaceData>}---extends
    the face topological-features of the Dcel with a record of type
    \cppLstinline{FaceData}.
    %
  \item \cppLstinline{Arr_extended_dcel<Traits, VertexData, HalfedgeData, FaceData, VertexBase, HalfedgeBase, FaceBase>}---extend
    all the topological-features of the Dcel.
  \end{compactenum}
  The \cppLstinline{VertexBase} template parameter must be substituted with
  either the base vertex type in the default Dcel, or an extended vertex type.
  Similarly, the \cppLstinline{HalfedgeBase} parameter must be substituted with
  the base halfedge type in the default Dcel, or an extended halfedge type,
  and the \cppLstinline{FaceBase} parameter must be substituted with the base
  face type in the default Dcel, or an extended face type.

  It is impossible to define a custom C++ type from Python code.
  Therefore, when the bindings are generated each one of the
  \cppLstinline{VertexData}, \cppLstinline{HalfedgeData}, and
  \cppLstinline{FaceData} template parameters must be substituted
  with a C++ type known at the time the bindings were implemented.
  Flexibility is nevertheless retained by substituting every one
  of these parameters with the generic Python object
  \cppLstinline{bp::object}. This object provides a
  generalized interface to Python objects.

  The \cmake{} flag \myLstinline{AOS2\_DCEL\_NAME} specifies which
  DCEL extension should be used for the generated bindings; see
  Table~\ref{tab:aos:dcel}

\end{itemize}

%%%%%%%% Arrangement Geometry Traits
\begin{table}[!htp]
  \caption{2D arrangement geometry traits options}
  \centering\begin{tikzpicture}
   \matrix (first) [tableName,
     column 1/.style={nodes={text width=18em}},
     column 2/.style={nodes={text width=31em}}
   ] {
     \myLstinline{AOS2\_GEOMETRY\_TRAITS\_NAME} & Type\\
     \myLstinline{nonCachingSegment} &
       \cppLstinline{Arr\_non\_caching\_segment\_basic\_traits\_2<Kernel>}\\
     \myLstinline{segment} & \cppLstinline{Arr\_segment_traits\_2<Kernel>}\\
     \myLstinline{linear} & \cppLstinline{Arr\_linear\_traits\_2<Kernel>}\\
     \myLstinline{conic} &
       \cppLstinline{Arr\_conic\_traits\_2<RatKernel, AlgKernel, NtTraits>}\\
     \myLstinline{circleSegment} &
       \cppLstinline{Arr\_circle\_segment_traits\_2<Kernel>}\\
     \myLstinline{algebraic} &
       \cppLstinline{Arr\_algebraic\_segment\_traits\_2<Coefficient>}\\
   };
  \end{tikzpicture}
  \label{tab:aos:gt}
\end{table}

%%%%%%%% Arrangement DCEL
\begin{table}[!htp]
 \caption{2D arrangement DCEL options}
  \centering\begin{tikzpicture}
  \matrix (first) [tableName,
     column 1/.style={nodes={text width=12em}},
     column 2/.style={nodes={text width=37em}}
   ] {
     \myLstinline{AOS2\_DCEL\_NAME} & Type\\
     \myLstinline{plain} & \cppLstinline{Arr_default_dcel<GeometryTraits>}\\
     \myLstinline{vertexExtended} &
       \cppLstinline{Arr_vertex_extended_dcel<Traits, Data>}\\
     \myLstinline{halfedgeExtended} &
       \cppLstinline{Arr_halfedge_extended_dcel<Traits, Data>}\\
     \myLstinline{faceExtended} &
       \cppLstinline{Arr_face_extended_dcel<Traits, Data>}\\
     \myLstinline{allExtended} &
       \cppLstinline{Arr_extended_dcel<Traits, VertexData, HalfedgeData, FaceData>}\\
  };
 \end{tikzpicture}
 \label{tab:aos:dcel}
\end{table}

The following code excerpt in the \cmake{} language sets the \cmake{} flags for our next example.
\framedInputCmake[\normalsize]{../../cmake/tests/aos2_epec_face_extended_release.cmake}

The following example constructs an arrangement with two faces. The
arrangement is induced by line segments and its face type is extended
with an integral label. The unbounded face and the other faced are labeled
'1' and '0', respectively.
\framedInputPython[\normalsize][15]{../../src/python_scripts/cgalpy_examples/aos2_fex.py}

% -----------------------------------------------------------------------------
\subsection{3D Triangulation Bindings}
\label{ssec:modules:tri3}
% -----------------------------------------------------------------------------
A triangulation of a set of points $\calP$ in \Rthree is a partition
of the convex hull of $\calP$ into tetrahedra whose vertices are the
points of $\calP$. Together with the unbounded cell having the convex
hull boundary as its frontier, the triangulation forms a partition
of \Rthree. Any two cells (3-face) are either disjoint or share a
common facet (2-face), edge (1-face) or vertex (0-face).
The binding module \iiiDTriangulationsPackage{} consists of bindings
of types provided by the \cgal{} packages \iiiDTriangulationsPackage{}
and \iiiDPeriodicTriangulationsPackage{}. Table~\ref{tab:tri3} lists
the \cmake{} flags associated with the \iiiDTriangulationsPackage{}
module. The default selection of the traits depends on the
triangulation selection; see Table~\ref{tab:tri3:traits-default}.

%%%%%%%% 3D Triangulation Module Flags
\begin{table}[!htp]
  \caption{\iiiDTriangulationsPackage{} module flags. For the default
  value of the traits name see Table~\ref{tab:tri3:traits-default}.}
  \centering\begin{tikzpicture}
  \matrix (first) [tableModuleFlags] {
    Name & Type & Default & Description\\
\myLstinline{TRI3\_NAME}                   & String & \myLstinline{regular}    & The 3D triangulation type\\
\myLstinline{TRI3\_VERTEX\_BASE\_NAME}     & String & \myLstinline{plain}      & The base vertex type\\
\myLstinline{TRI3\_CELL\_BASE\_NAME}       & String & \myLstinline{plain}      & The base cell type\\
\myLstinline{TRI3\_TRAITS\_NAME}           & String & *                        & The geometry traits of a 3D triangulation\\
\myLstinline{TRI3\_CONCURRENCY\_NAME}      & String & \myLstinline{sequential} & The concurrency method\\
\myLstinline{TRI3\_LOCATION\_POLICY\_NAME} & String & \myLstinline{compact}    & The location policy\\
  };
  \end{tikzpicture}
  \label{tab:tri3}
\end{table}

%%%%%%%% 3D triangulation Traits default selection
\begin{table}
  \caption{3D triangulation Traits default selection}
  \centering\begin{tikzpicture}
  \matrix (first) [tableName,
    column 1/.style={nodes={text width=30em}},
    column 2/.style={nodes={text width=19em}}
  ] {
    3D Triangulation Type & Default \myLstinline{TRI3\_TRAITS\_NAME}\\
    \cppLstinline{Periodic\_3\_triangulation\_3<>} &
      \myLstinline{periodicPlain}\\
    \cppLstinline{Periodic\_3\_regular_triangulation\_3<>} &
      \myLstinline{periodicRegular}\\
    \cppLstinline{Periodic\_3\_Delaunay_triangulation\_3<>} &
      \myLstinline{periodicDelaunay}\\
    \textit{Otherwise} & \myLstinline{kernel}\\
  };
  \end{tikzpicture}
  \label{tab:tri3:traits-default}
\end{table}

The \iiiDTriangulationsPackage{} package~\cite{cgal:pt-t3-21a}
provides several class templates, instances of which can be used to
represent a variaty 3D triangulations. In particular, the package
provides the following class templates:
\begin{compactenum}
\item \cppLstinline{Triangulation\_3<Traits\_3, DataStructure\_3, SpatialLockDataStructure>},
\item \cppLstinline{Regular\_triangulation\_3<Traits\_3, DataStructure\_3>}, and
\item \cppLstinline{Delaunay\_triangulation\_3<Traits\_3, DataStructure\_3, LocationPolicy>}
\end{compactenum}
Instance of the \cppLstinline{Delaunay\_triangulation\_3<>} and
\cppLstinline{Regular\_triangulation\_3<>} class templates can be used
to represent Delaunay and regular triangulation, respectively. In a
regular triangulation points have an associated weight, and some
points can be hidden and do not result in vertices in the
triangulation.

The \iiiDPeriodicTriangulationsPackage{}
package~\cite{cgal:ct-pt3-21a} supports triangulations of sets of
points in the three-dimensional flat torus. This package provides the
class templates
\begin{compactenum}
\item \cppLstinline{Periodic\_3\_triangulation\_3<Traits\_3, DataStructure\_3>},
\item \cppLstinline{Periodic\_3\_regular\_triangulation\_3<Traits\_3, DataStructure\_3>}, and
\item \cppLstinline{Periodic\_3\_Delaunay\_triangulation\_3<Traits\_3, DataStructure\_3>}
\end{compactenum}

The \cmake{} flag \myLstinline{TRI3\_NAME} specifies the particular
types of triangulations, bindings for which should be generated; see
Table~\ref{tab:tri3:names}.

\begin{table}[!htp]
  \caption{3D triangulation name options}
  \centering\begin{tikzpicture}
  \matrix (first) [tableName,
    column 1/.style={nodes={text width=10em}},
    column 2/.style={nodes={text width=39em}}
  ] {
    \myLstinline{TRI3\_NAME} & Type\\
    \myLstinline{plain} & \cppLstinline{Triangulation\_3<Traits\_3, DataStructure\_3>}\\
    \myLstinline{regular} & \cppLstinline{Regular\_triangulation\_3<Traits\_3, DataStructure\_3>}\\
    \myLstinline{delaunay} & \cppLstinline{Delaunay\_triangulation\_3<Traits\_3, DataStructure\_3>}\\
    \myLstinline{periodicPlain} & \cppLstinline{Periodic\_3\_triangulation\_3<Traits\_3, DataStructure\_3>}\\
    \myLstinline{periodicRegular} & \cppLstinline{Periodic\_3\_regular\_triangulation\_3<Traits\_2, DataStructure\_3>}\\
    \myLstinline{periodicDelaunay} & \cppLstinline{Periodic\_3\_Delaunay\_triangulation\_3<Traits\_2, DataStructure\_3>}\\
  };
  \end{tikzpicture}
  \label{tab:tri3:names}
\end{table}

When any template above is instantiated the template parameter
\cppLstinline{Traits\_3} must be substituted with a model of a
respective geometric traits concept,
\cppLstinline{TriangulationTraits\_3},
\cppLstinline{RegularTriangulationTraits\_3},
\cppLstinline{DelaunayTriangulationTraits\_3},
\cppLstinline{Periodic\_3TriangulationTraits\_3},
\cppLstinline{Periodic\_3RegularTriangulationTraits\_3}, or
\cppLstinline{Periodic\_3DelaunayTriangulationTraits\_3}. The geometric
traits class provides the type of points to use as well as elementary
operations on them. Observe that the Kernel is a model of the former
three concepts and the only choice for the binding of instances of
\cppLstinline{Triangulation\_3<>},
\cppLstinline{Regular\_triangulation\_3<>}, or
\cppLstinline{Delaunay\_triangulation\_3<>} (the non-periodic
triangulation types). The \cmake{} flag
\myLstinline{TRI3\_TRAITS\_NAME} specifies the type of traits to be
used for the generated bindings; see table~\ref{tab:tri3:traits}.

\begin{table}[!htp]
  \caption{3D triangulation traits options}
  \centering\begin{tikzpicture}
  \matrix (first) [tableName,
    column 1/.style={nodes={text width=12em}},
    column 2/.style={nodes={text width=37em}}
  ] {
    \myLstinline{TRI3\_TRAITS\_NAME} & Type\\
    \myLstinline{kernel} & Kernel \\
    \myLstinline{periodicPlain} &
      \cppLstinline{Periodic\_3\_triangulation\_traits\_3<Kernel, Offset>}\\
    \myLstinline{periodicRegular} &
      \cppLstinline{Periodic\_3\_regular\_triangulation\_traits\_3<Kernel, Offset>}\\
    \myLstinline{periodicDelaunay} &
      \cppLstinline{Periodic\_3\_Delaunay\_triangulation\_traits\_3<Kernel, Offset>}\\
  };
  \end{tikzpicture}
  \label{tab:tri3:traits}
\end{table}

The type that substitutes the \cppLstinline{DataStructure\_3}
parameter models the concept
\cppLstinline{TriangulationDataStructure\_3}. An object of this type
stores the combinatorial structure of the triangulation; it is an
instance of the class template
\cppLstinline{Triangulation\_data\_structure\_3<VertexBase, CellBase,
  ConcurrencyTag>}.

The types that substitute the \cppLstinline{VertexBase} and
\cppLstinline{CellBase} template parameters when the template
\cppLstinline{Triangulation\_data\_structure\_3<>} is instantiated
must model the concepts \cppLstinline{TriangulationDSVertexBase\_3}
and \cppLstinline{TriangulationDSCellBase\_3}, respectively, when the
triangulation instance is not used to define an alpha shape type. If
the triangulation is used to define a plain alpha shape type, these
parameters must be substituted with types that model the concepts
\cppLstinline{AlphaShapeVertex\_3} and
\cppLstinline{AlphaShapeCell\_3}, respectively; see
Section~\ref{ssec:modules:as3}. If the triangulation is used to define
a fixed alpha shape type, these parameters must be substituted with
types that model the concepts \cppLstinline{AlphaShapeVertex\_3} and
\cppLstinline{AlphaShapeCell\_3}, respectively.  These types represent
the base type of the vertex and the base type of the cell of the
triangulation, respectively. The \cmake{} flags
\myLstinline{TRI3\_VERTEX\_BASE\_NAME} and
\myLstinline{TRI3\_CELL\_BASE\_NAME} specify the vertex and base
types, respectively; see Table~\ref{tab:tri3:vertex}.

\begin{table}[!htp]
  \caption{3D triangulation vertex and cell base options.
    \cppLstinline{Tri...} stands for \cppLstinline{Triangulation};
    \cppLstinline{entity} must be replaced with either
    \cppLstinline{vertex} or \cppLstinline{cell};
    \cppLstinline{K} and \cppLstinline{Size} are aliases for
    \cppLstinline{Kernel} and \cppLstinline{size\_t}, respectively;
    \cppLstinline{Ec} is either \cppLstinline{true} or
    \cppLstinline{false} based on the selection of
    \cmakeLstinline{AS3\_EXACT\_COMPARISON}; see Table~\ref{tab:as}.
  }
  \centering\begin{tikzpicture}
  \matrix (first) [tableName,
    row 1/.append style={nodes={minimum height=2.8em}},
    column 1/.style={nodes={text width=18.1em}},
    column 2/.style={nodes={text width=30.9em}}
  ] {%
    \setlength{\tabcolsep}{0pt}\renewcommand{\arraystretch}{0}
    \begin{tabular}{c}
      \cmakeLstinline{TRI3\_VERTEX\_BASE\_NAME}
      \textcolor{white}{or}\\[5pt]
      \textcolor{white}{\cmakeLstinline{TRI3\_CELL\_BASE\_NAME}}
    \end{tabular} & Type\\
    \myLstinline{plain} &
      \cppLstinline{Vb = Tri...\_entity\_base\_3<K>}\\
    \myLstinline{plainWithInfo} &
      \cppLstinline{Vbi = Tri...\_entity\_base\_with\_info\_3<Size, K>}\\
    \myLstinline{regular} &
      \cppLstinline{Rvb = Regular\_triangulation\_entity\_base\_3<K>}\\
    \myLstinline{regualrWithInfo} &
      \cppLstinline{Rvbi = Tri...\_entity\_base\_with\_info\_3<Size, K, Rvb>}\\
    \myLstinline{alphaShape} &
      \cppLstinline{Alpha\_shape\_entity\_base\_3<K, Vb, Ec>}\\
    \myLstinline{alphaShapeWithInfo} &
      \cppLstinline{Alpha\_shape\_entity\_base\_3<K, Vbi, Ec>}\\
    \myLstinline{alphaShapeRegular} &
      \cppLstinline{Alpha\_shape\_entity\_base\_3<K, Rvb, Ec>}\\
    \myLstinline{alphaShapeRegularWithInfo} &
      \cppLstinline{Alpha\_shape\_entity\_base\_3<K, Rvbi, Ec>}\\
    \myLstinline{fixedAlphaShape} &
      \cppLstinline{Fixed\_alpha\_shape\_entity\_base\_3<K, Vb>}\\
    \myLstinline{fixedAlphaShapeWithInfo} &
      \cppLstinline{Fixed\_alpha\_shape\_entity\_base\_3<K, Vbi>}\\
    \myLstinline{fixedAlphaShapeRegular} &
      \cppLstinline{Fixed\_alpha\_shape\_entity\_base\_3<K, Rvb>}\\
    \myLstinline{fixedAlphaShapeRegularWithInfo} &
      \cppLstinline{Fixed\_alpha\_shape\_entity\_base\_3<K, Rvbi>}\\
  };
  \end{tikzpicture}
  \label{tab:tri3:vertex}
\end{table}

The template parameter \cppLstinline{ConcurrencyTag} is substituted
with either \cppLstinline{Sequential\_tag} or
\cppLstinline{Parallel\_tag} when the template
\cppLstinline{Triangulation\_data\_structure\_3<VertexBase, CellBase,
  ConcurrencyTag} is instantiated. It enables the use of a concurrent
container to store vertices and cells. The \cmake{} flag
\cmakeLstinline{TRI3\_CONCURRENCY\_NAME} determines the selection; see
Table~\ref{tab:tri3:concurrency}.

\begin{table}[!htp]
  \caption{3D triangulation concurrency options}
  \centering\begin{tikzpicture}
  \matrix (first) [tableName,
    column 1/.style={nodes={text width=18em}},
    column 2/.style={nodes={text width=31em}}
  ] {
    \cmakeLstinline{TRI3\_CONCURRENCY\_NAME} & Type\\
    \myLstinline{sequential} & \cppLstinline{Sequential\_tag}\\
    \myLstinline{parallel} & \cppLstinline{Parallel\_tag}\\
  };
  \end{tikzpicture}
  \label{tab:tri3:concurrency}
\end{table}
%

The template parameter \cppLstinline{ConcurrencyTag} is substituted
with either \cppLstinline{Fast\_location} or
\cppLstinline{Compact\_location} when the template
\cppLstinline{Delaunay\_triangulation\_3<Traits\_3, DataStructure\_3,
  LocationPolicy>} is instantiated. It enables a faster point location
at the account of memory space. This is useful when performing point
locations or random point insertions in large data sets.  The \cmake{}
flag \cmakeLstinline{TRI3\_LOCATION\_POLICY\_NAME} determines the
selection; see Table~\ref{tab:tri3:location}.

\begin{table}[!htp]
  \caption{3D triangulation location policy options}
  \centering\begin{tikzpicture} \matrix (first) [tableName, column
    1/.style={nodes={text width=18em}}, column 2/.style={nodes={text
        width=31em}} ] { \cmakeLstinline{TRI3\_LOCATION\_POLICY\_NAME}
    & Type\\ \myLstinline{fast} &
    \cppLstinline{Fast\_location}\\ \myLstinline{compact} &
    \cppLstinline{Compact\_location}\\ };
  \end{tikzpicture}
  \label{tab:tri3:location}
\end{table}

The following code excerpt in the \cmake{} language sets the \cmake{}
flags for our next example.
\framedInputCmake[\normalsize]{../../cmake/tests/tri3_del_epic_release.cmake}

The following example constructs a three-dimensional Delaunay
triangulation from six points and verivies that the triangulation is valid.
\framedInputPython[\normalsize][15]{../../src/python_scripts/cgalpy_examples/tri3_del.py}

% -----------------------------------------------------------------------------
\subsection{3D Alpha Shapes}
\label{ssec:modules:as3}
% -----------------------------------------------------------------------------
Given a set of points sampled in a 3D body, an alpha shape
(a.k.a.\ alpha complex) is demarcated by a frontier, which is a linear
approximation of the original boundary of the body. An alpha shape
object maintains an underlying triangulation of a set of input points.
There are two distinguished versions of alpha shapes as follows. Basic
alpha shapes are based on the Delaunay triangulation and weighted
alpha shapes are based on its generalization, the regular
triangulation, where the euclidean distance is replaced by the power
to weighted points. In a plain alpha shape each $k$-face of this
triangulation is associated with an interval specifying for which
values of $\alpha$ the face belongs to the alpha complex. In a fixed
alpha shape each $k$-face is associated with a classification that
specifies its status in the alpha complex, alpha being fixed. The
package \iiDArrangementsPackage{}~\cite{cgal:dy-as3-21a} provides the
class templates \cppLstinline{Alpha\_shape\_3<Dt,
  ExactAlphaComparisonTag>} and
\cppLstinline{Fixed\_alpha\_shape\_3<Dt>}. Instances of both templates
can be used to represent a large variaty of alpha shapes for a given
set of points. Table~\ref{tab:as3} lists the \cmake{} flags associated
with the \iiDArrangementsPackage{}.

%%%%%%%% 3D Alpha Shape flags
\begin{table}[!htp]
  \caption{\iiiDAlphaShapesPackage{} module flags}
  \centering\begin{tikzpicture}
  \matrix (first) [tableModuleFlags,
    row 3/.style={nodes={minimum height=2.8em}}
  ] {
Name & Type &  Default & Description\\
\myLstinline{AS3\_NAME}              & String  & \myLstinline{plain} & The 3D Alpha shape type\\
\myLstinline{AS3\_EXACT\_COMPARISON} & Boolean & \myLstinline{false} & Determines whether to apply exact comparisons\\
  };
 \end{tikzpicture}
 \label{tab:as3}
\end{table}

The \cmake{} flag \myLstinline{AS3\_NAME} specifies which alpha shape
should be used for the generated bindings; see
Table~\ref{tab:as3:types}

%%%%%%%% 3D Alpha Shapes Types
\begin{table}[!htp]
 \caption{\iiiDAlphaShapesPackage{} types}
  \centering\begin{tikzpicture}
  \matrix (first) [tableName,
     column 1/.style={nodes={text width=12em}},
     column 2/.style={nodes={text width=37em}}
   ] {
     \myLstinline{AS3\_NAME} & Type\\
     \myLstinline{plain} & \cppLstinline{Alpha\_shape\_3<>}\\
     \myLstinline{fixed} & \cppLstinline{Fixed\_alpha\_shape\_3<>}\\
  };
 \end{tikzpicture}
 \label{tab:as3:types}
\end{table}

The template parameter
\cppLstinline{Dt} must be substituted with an instance of one of the
class templates that can be used to represent a triangulation, namly:
\begin{compactenum}
\item \cppLstinline{Delaunay_triangulation_3},
\item \cppLstinline{Regular_triangulation_3},
\item \cppLstinline{Periodic_3_Delaunay_triangulation_3}, or
\item \cppLstinline{Periodic_3_regular_triangulation_3};
\end{compactenum}
see Section~\ref{ssec:modules:tri3}.
Note that \cppLstinline{Dt::Geom\_traits}, \cppLstinline{Dt::Vertex}, and
\cppLstinline{Dt::Face} must be models the concepts
\cppLstinline{AlphaShapeTraits\_3}, \cppLstinline{AlphaShapeVertex\_3},
and \cppLstinline{AlphaShapeCell\_3}, respectively.

% -----------------------------------------------------------------------------
\subsection{Spatial Searching}
\label{ssec:modules:ss}
% -----------------------------------------------------------------------------
The \dDSpatialSearchingPackage{} package~\cite{cgal:tf-ssd-21a}
implements exact and approximate distance browsing by providing
implementations of algorithms supporting
\begin{compactenum}
\item both nearest and furthest neighbor searching,
\item both exact and approximate searching,
\item (approximate) range searching,
\item (approximate) k-nearest and k-furthest neighbor searching,
\item (approximate) incremental nearest and incremental furthest neighbor searching, and
\item query items representing points and spatial objects.
\end{compactenum}
In these searching problems a set $\calP$ of data points in
$d$-dimensional space is given. These points in $\calP$ are
preprocessed into a tree data structure, so that given any query item
$q$ the points of $\calP$ can be browsed efficiently. The approximate
\dDSpatialSearchingPackage{} package is designed for data sets that
are small enough to store the search structure in main memory (in
contrast to approaches from databases that assume that the data
reside in secondary storage).

%%%%%%%%
\begin{table}[!htp]
  \caption{\dDSpatialSearchingPackage{} module flags}
  \centering\begin{tikzpicture}
  \matrix (first) [tableModuleFlags,
    row 2/.style={nodes={minimum height=2.8em}}
  ] {
Name & Type &  Default & Description\\
\myLstinline{SPATIAL\_SEARCHING\_DIMENSION} & Integer & \myLstinline{2} & The dimension of spatial searching related classes\\
  };
 \end{tikzpicture}
 \label{tab:ss}
\end{table}
